\documentclass[11pt,a4paper]{article}

\usepackage[utf8]{inputenc}
\usepackage[T1]{fontenc}
\usepackage[english]{babel}
\usepackage{graphicx}
\usepackage{caption}
\usepackage{subcaption}
\usepackage{amsmath}
\usepackage{amssymb}
\usepackage{hyperref}
\usepackage{csquotes}
\usepackage{geometry}
\usepackage{fancyhdr}
\usepackage{booktabs}
\usepackage{longtable}
\usepackage{enumitem}
\usepackage{url}

\geometry{left=2.5cm, right=2.5cm, top=2.5cm, bottom=2.5cm}
\setlength{\parindent}{0pt}
\setlength{\parskip}{6pt}

\hypersetup{
    colorlinks=true,
    linkcolor=blue,
    urlcolor=blue,
    citecolor=blue
}

\newcommand{\VLM}{\textsc{VLM}}

\pagestyle{fancy}
\fancyhf{}
\rhead{CVCS 2026 - Do VLMs Really See?}
\lhead{Hallucination Evaluation}
\rfoot{\thepage}

\title{Do Vision-Language Models Really See What They Say?}
\author{MDN-C-CS Team}
\date{May 2026}

\begin{document}

\maketitle

\begin{abstract}
This project evaluates whether Vision-Language Models (VLMs) hallucinate visual content that is not supported by an image. We focus on four hallucination categories: absent objects, incorrect attributes, spatial relations, and counting errors. The implementation provides a synthetic probing benchmark, evaluates LLaVA-v1.5 and Qwen3-VL by default, reports POPE-style false-positive metrics, and tests a prompt-engineering mitigation strategy.
\end{abstract}

\section{Introduction}

Vision-Language Models can answer questions about images and generate fluent descriptions, but they may also hallucinate: they produce plausible statements that are not grounded in the visual evidence. This failure mode is important for safety-critical applications such as medical imaging, autonomous driving, and media verification.

\subsection{Hallucination Types}

\begin{itemize}
    \item \textbf{Object hallucination}: the model claims that an absent object is present.
    \item \textbf{Attribute hallucination}: the model assigns an incorrect color or property.
    \item \textbf{Relation hallucination}: the model misreads spatial relationships.
    \item \textbf{Counting hallucination}: the model reports an incorrect number of objects.
\end{itemize}

\subsection{Contributions}

\begin{enumerate}
    \item A self-consistent probing benchmark for four VLM hallucination categories.
    \item A comparative evaluation setup for LLaVA-v1.5 and Qwen3-VL, with optional InternVL3.5 support.
    \item POPE-style metrics and condition-level error analysis.
    \item A baseline vs strict-prompt mitigation study.
\end{enumerate}

\section{Evaluated Models}

\subsection{LLaVA-v1.5}

LLaVA-v1.5 combines a visual encoder with a Vicuna-based language model and is a widely used open VLM baseline.

\subsection{Qwen3-VL}

Qwen3-VL is a recent multimodal model family from the Qwen series. The project uses the 8B Instruct checkpoint by default.

\subsection{Optional: InternVL3.5}

InternVL3.5 is included as an optional third model when additional GPU time is available.

\section{Methodology}

\subsection{Benchmark Design}

The benchmark consists of image-question pairs where the correct answer is always \enquote{no}. This makes it a negative probing benchmark similar in spirit to POPE: a \enquote{yes} answer is a false-positive hallucination.

\begin{table}[h]
\centering
\caption{Trap-question categories}
\begin{tabular}{@{}lll@{}}
\toprule
\textbf{Type} & \textbf{Description} & \textbf{Example} \\
\midrule
Object absent & Query about an absent object & \enquote{Is there a dog in this image?} \\
Attribute & Query with a false property & \enquote{Is the square red?} \\
Relation & Query with a false spatial relation & \enquote{Is the circle above the triangle?} \\
Count & Query with a false count & \enquote{Are there five circles?} \\
\bottomrule
\end{tabular}
\end{table}

\subsection{Metrics}

We report:

\begin{itemize}
    \item Overall accuracy.
    \item Hallucination rate: fraction of \enquote{yes} answers on negative probes.
    \item POPE-style false-positive rate.
    \item Accuracy by question type.
    \item Condition-level metrics by object, claimed color, relation, and count.
\end{itemize}

\subsection{Mitigation}

The baseline prompt asks the model to answer only yes or no. The mitigation prompt adds an explicit visual-verification instruction:

\begin{quote}
Answer only yes or no. Verify the visual evidence carefully. If the requested object, attribute, relation, or count is not clearly visible, answer no. Do not infer likely objects from context.
\end{quote}

\section{Experiments and Results}

Replace the placeholder values below with the generated JSON metrics after running the cluster jobs.

\begin{table}[h]
\centering
\caption{Model-level results}
\begin{tabular}{@{}lccc@{}}
\toprule
\textbf{Model} & \textbf{Accuracy} & \textbf{Hallucination Rate} & \textbf{Unknown Rate} \\
\midrule
LLaVA-v1.5 & -- & -- & -- \\
Qwen3-VL & -- & -- & -- \\
\bottomrule
\end{tabular}
\end{table}

\begin{figure}[h]
\centering
\includegraphics[width=0.8\textwidth]{reports/llava15/complete_results.png}
\caption{Example generated plot for LLaVA-v1.5. Update the path as needed.}
\end{figure}

\section{Analysis}

Discuss which question categories and metadata conditions produce the highest hallucination rates. In particular, compare object-absent questions against relation and count questions, because relation and count probes often require more precise visual grounding.

\section{Mitigation Results}

Compare baseline folders such as \texttt{results/llava15_baseline} with strict-prompt folders such as \texttt{results/llava15_strict}. Report whether the strict prompt reduces false-positive hallucination and whether it increases unknown or overly conservative answers.

\section{Limitations}

The default benchmark is synthetic. For a stronger scientific evaluation, add real-image probes from MSCOCO, Visual Genome, MMHal-Bench, or HallusionBench and compare whether the same error patterns hold.

\section{Conclusion}

This project provides a reproducible workflow for probing VLM hallucination, comparing recent open models, analyzing failure conditions, and testing a prompt-based mitigation strategy.

\section*{References}

\begin{itemize}
    \item Rohrbach et al. Object Hallucination in Image Captioning. EMNLP 2018.
    \item Li et al. Evaluating Object Hallucination in Large Vision-Language Models. EMNLP 2023.
    \item Liu et al. Improved Baselines with Visual Instruction Tuning (LLaVA-v1.5). CVPR 2024.
    \item Bai et al. Qwen3-VL Technical Report. arXiv 2025.
    \item Wang et al. InternVL3.5: Advancing Open-Source Multimodal Models in Versatility, Reasoning, and Efficiency. arXiv 2025.
    \item Guan et al. HallusionBench: An Advanced Diagnostic Suite for Entangled Language Hallucination and Visual Illusion in Large Vision-Language Models. CVPR 2024.
    \item Compagnoni et al. Mitigating Hallucinations in Multimodal LLMs via Object-aware Preference Optimization. BMVC 2025.
\end{itemize}

\end{document}
